\documentclass[8pt,a4paper]{extarticle}
\usepackage[margin=0.4in]{geometry}
\usepackage{amsmath,amssymb,amsthm}
\usepackage{multicol}
\usepackage{enumitem}

\setlength{\parindent}{0pt}
\setlength{\parskip}{2pt}
\setlength{\columnsep}{16pt}

\newcommand{\Z}{\mathbb{Z}}
\newcommand{\Q}{\mathbb{Q}}
\newcommand{\F}{\mathbb{F}}
\newcommand{\lcm}{\operatorname{lcm}}
\newcommand{\chr}{\operatorname{char}}

% Section styling
\makeatletter
\renewcommand\section{\@startsection{section}{1}{\z@}%
  {-2ex plus -0.5ex minus -0.2ex}%
  {0.8ex plus 0.2ex}%
  {\normalfont\large\bfseries}}
\renewcommand\subsection{\@startsection{subsection}{2}{\z@}%
  {-1.5ex plus -0.5ex minus -0.2ex}%
  {0.5ex plus 0.1ex}%
  {\normalfont\normalsize\bfseries}}
\makeatother

\begin{document}

%=== TITLE PAGE ===%
\begin{titlepage}
	\centering
	\makeatletter
	\renewcommand{\Huge}{\@setfontsize\Huge{26pt}{30pt}}
	\renewcommand{\LARGE}{\@setfontsize\LARGE{20pt}{24pt}}
	\renewcommand{\Large}{\@setfontsize\Large{16pt}{20pt}}
	\renewcommand{\large}{\@setfontsize\large{13pt}{16pt}}
	\makeatother

	\vspace*{1cm}
	{\Huge \textbf{Algebra Reference Sheet}} \\
	\vspace{20pt}
	{\LARGE MATH-310 — Anna Lachowska} \\
	\vspace*{1cm}
	{\Large Notes by Ali EL AZDI} \\
	\vfill
  \textbf{Abstract.} This sheet summarizes the main techniques and theorems needed to solve algebra exams. Good luck!
\vfill
	{\large January 2026}
	\vspace*{60px}
\end{titlepage}

%=== PAGE 1: THEORY ===%
\begin{multicols}{2}
\raggedcolumns

\section*{Part I: Group Theory}

\subsection*{1. Basic Definitions}

\textbf{Group.} Set $G$ with operation satisfying: closure, associativity, identity $e$, inverses.

\textbf{Order.} $|G|$ = number of elements. $o(g)$ = smallest $n > 0$ with $g^n = e$.

\textbf{Subgroup.} $H \le G$: nonempty subset closed under operation and inverses.

\textbf{Normal subgroup.} $H \triangleleft G$ iff $gHg^{-1} = H$ for all $g \in G$.
\begin{itemize}[leftmargin=*, itemsep=0pt, topsep=2pt]
\item Index 2 subgroups are always normal.
\item To check: verify $gHg^{-1} = H$ only for generators of $G$.
\end{itemize}

\subsection*{2. Group Actions \& Subgroups}

\textbf{Orbit.} $\text{Orb}_G(x) = G \cdot x = \{g \cdot x : g \in G\}$

\textbf{Stabilizer.} $\text{Stab}_G(x) = G_x = \{g \in G : g \cdot x = x\}$

\textbf{Orbit-Stabilizer.} $|G| = |\text{Orb}(x)| \cdot |\text{Stab}(x)|$

\textbf{Center.} $Z(G) = \{g \in G : gx = xg \ \forall x \in G\}$ \quad (always $\triangleleft G$; for $p$-groups: $\neq \{e\}$)

\textbf{Centralizer.} $C_G(x) = \{h \in G : hx = xh\}$

Note: $Z(G) \subseteq C_G(x)$ for all $x$. Also $C_G(x) = \text{Stab}(x)$ under conjugation action.

\textbf{Class equation.} $|G| = |Z(G)| + \sum_i [G : C_G(x_i)]$

where $x_i$ are representatives of non-trivial conjugacy classes.

\subsection*{3. Key Theorems}

\textbf{Lagrange.} $H \le G$ finite $\Rightarrow$ $|H|$ divides $|G|$.

\textbf{Cauchy.} $p$ prime, $p \mid |G|$ $\Rightarrow$ $\exists g \in G$ with $o(g) = p$.

\textbf{$|G| = p$ prime.} $G$ is cyclic (hence abelian).

\textbf{$|G| = p^2$.} $G$ is abelian.

\textit{Why?} Class equation $\Rightarrow$ $|Z(G)| \geq p$. If $|Z(G)| = p$, then $G/Z(G)$ cyclic $\Rightarrow$ $G$ abelian. So $Z(G) = G$.

\subsection*{4. Homomorphisms}

\textbf{Group hom.} $\phi: G \to H$ with $\phi(ab) = \phi(a)\phi(b)$.

\textbf{Ring hom.} $\phi: R \to S$ with $\phi(a+b) = \phi(a)+\phi(b)$, $\phi(ab) = \phi(a)\phi(b)$, $\phi(1) = 1$.

\textbf{Isomorphism.} Bijective homomorphism.

\textbf{Kernel.} $\ker(\phi) = \phi^{-1}(e)$. For groups: $\ker \triangleleft G$. For rings: $\ker$ is an ideal.

\textbf{First Isomorphism Thm.} $G/\ker(\phi) \cong \text{Im}(\phi)$.

\vspace{4pt}
\fbox{\parbox{\dimexpr\linewidth-2\fboxsep-2\fboxrule}{%
\textbf{Hom existence:}\\[2pt]
$\bullet$ Group: $\text{Hom}(C_m, C_n) \neq 0 \iff \gcd(m,n) > 1$.\\[1pt]
\textit{\small ($C_n$ has element of order $d \mid m$ iff $d \mid \gcd(m,n)$)}\\[2pt]
$\bullet$ Ring: $\Z/m\Z \to \Z/n\Z$ (unital) exists iff $n \mid m$.
}}
\vspace{4pt}

\subsection*{5. Symmetric Groups}

$S_n$ = bijections of $\{1,\ldots,n\}$, \ $|S_n| = n!$

\textbf{Cycle notation.} $(a_1 \cdots a_k)$: $a_i \mapsto a_{i+1}$, $a_k \mapsto a_1$.

\textbf{Order.} Disjoint cycles: $\lcm$ of lengths.

\textbf{Parity.} $k$-cycle = $(k-1)$ transpositions. Even iff $k$ odd.

\textbf{$A_n$.} Even permutations. $|A_n| = n!/2$. Contains: $e$, 3-cycles, double transpositions.

\textbf{Conjugation.} $\sigma(i_1 \cdots i_k)\sigma^{-1} = (\sigma(i_1) \cdots \sigma(i_k))$.

Preserves cycle type. Conjugacy classes = cycle types.

\vspace{4pt}
\fbox{\parbox{\dimexpr\linewidth-2\fboxsep-2\fboxrule}{%
\textbf{Non-disjoint products:}\\[2pt]
$(1\,2\cdots k)(i\,m)$, $i \le k < m$: becomes $(k{+}1)$-cycle.\\[2pt]
$(1\,2\cdots k)(i\,j)$, $i < j \le k$: splits into two cycles.
}}
\vspace{4pt}

\subsection*{6. Finite Abelian Groups}

\textbf{Structure Thm.} $G \cong \Z_{p_1^{e_1}} \times \cdots \times \Z_{p_k^{e_k}}$ (prime powers).

\textbf{Counting.} \# abelian groups of order $n = p_1^{a_1} \cdots p_r^{a_r}$ is $\prod_i p(a_i)$, where $p(a)$ = partitions of $a$.

\textit{Partitions of small $n$: $p(1)=1$, $p(2)=2$, $p(3)=3$, $p(4)=5$, $p(5)=7$.}

\textbf{Invariant factors.} $d_1 \mid d_2 \mid \cdots \mid d_r$ with $G \cong \Z_{d_1} \times \cdots \times \Z_{d_r}$.

\vspace{4pt}
\fbox{\parbox{\dimexpr\linewidth-2\fboxsep-2\fboxrule}{%
\textbf{Order existence in finite abelian group:}\\[2pt]
$\bullet$ Max element order = $d_r$ (largest invariant factor).\\[1pt]
$\bullet$ Element of order $n$ exists iff $n \mid d_r$.\\[2pt]
\textbf{In cyclic group $C_m$:} Element of order $n$ exists iff $n \mid m$.\\[1pt]
\textit{\small (Cyclic group of order $m$ has exactly one subgroup of order $d$ for each $d \mid m$.)}
}}
\vspace{4pt}

\textbf{Order in products.} For $(a,b) \in C_m \times C_n$: $\text{ord}(a,b) = \lcm(\text{ord}(a), \text{ord}(b))$.

Max order in $C_m \times C_n$ is $\lcm(m,n)$. Generalizes: max order in $C_{n_1} \times \cdots \times C_{n_k}$ is $\lcm(n_1, \ldots, n_k)$.

\textbf{Cyclic test.} $\Z_m \times \Z_n \cong \Z_{mn}$ iff $\gcd(m,n) = 1$.

\textit{Why?} $\lcm(m,n) = mn$ iff $\gcd(m,n) = 1$. Only then $\exists$ element of order $mn$.

\section*{Part II: Rings \& Fields}

\subsection*{7. Ring Hierarchy}

Field $\subset$ Euclidean Domain $\subset$ PID $\subset$ Integral Domain $\subset$ Ring

\textbf{Integral domain.} Commutative, $1 \neq 0$, no zero divisors.

\textbf{Zero divisor.} $a \neq 0$ with $ab = 0$ for some $b \neq 0$.

\textbf{Unit.} Has multiplicative inverse. Units are never zero divisors.

\textbf{Irreducible.} $c \neq 0$, not unit, $c = ab \Rightarrow a$ or $b$ unit.

\subsection*{8. Fields \& Characteristic}

\textbf{Field.} Every nonzero element is a unit.

\textbf{Finite fields.} $|\F_q| = p^n$ for prime $p$. Unique up to isomorphism.

$\F_q$ exists iff $q = p^n$. (No $\F_{12}$, $\F_{18}$.)

\textbf{Unit group.} $\F_q^\times$ is \textbf{cyclic} of order $q-1$.

\textbf{Characteristic.} Let $\tau: \Z \to R$ be the unique ring hom ($\tau(1) = 1_R$).

$\chr(R) = n$ where $\ker(\tau) = n\Z$. Equivalently: $\chr(R) = \min\{n > 0 : n \cdot 1_R = 0\}$, or $0$ if none.

\vspace{4pt}
\fbox{\parbox{\dimexpr\linewidth-2\fboxsep-2\fboxrule}{%
$\chr(\Z/n\Z) = n$ \quad $\chr(\F_p) = p$ \quad $\chr(\Q) = 0$\\[2pt]
$\chr(A \times B) = \lcm(\chr A, \chr B)$
}}
\vspace{4pt}

\subsection*{9. Ideals \& Quotients}

\textbf{Ideal operations} (in PID):
\begin{itemize}[leftmargin=*, itemsep=0pt, topsep=2pt]
\item $(a)(b) = (ab)$
\item $(a) \cap (b) = (\lcm(a,b))$
\item $(a) + (b) = (\gcd(a,b))$
\end{itemize}

\textit{Ex: $(6\Z) \cap (15\Z) \cap (10\Z) = (30\Z)$.}

\textbf{Quotient ring.} $R/I$: elements are cosets $[a] = a + I$.

\textbf{In $F[X]/(p)$:}
\begin{itemize}[leftmargin=*, itemsep=0pt, topsep=2pt]
\item Size: $|F|^{\deg p}$. Elements: polynomials of degree $< \deg p$.
\item Field iff $p$ irreducible. (Irreducible $\Rightarrow$ $(p)$ maximal.)
\item $\bar{f}$ unit iff $\gcd(f,p) = 1$. \ $\bar{f}$ zero divisor iff $\gcd(f,p) \neq 1$.
\end{itemize}

\subsection*{10. Irreducibility Tests}

\textbf{Degree $\le 3$.} Irreducible iff no roots in $F$.

\textbf{Rational root test.} $f \in \Z[X]$, root $\frac{r}{s}$ reduced $\Rightarrow$ $r \mid a_0$, $s \mid a_n$.

\textbf{Eisenstein.} Prime $p$: $p \nmid a_n$, $p \mid a_i$ for $i < n$, $p^2 \nmid a_0$ $\Rightarrow$ irred.\ over $\Q$.

\textbf{Reduction mod $p$.} $\bar{f} \in \F_p[X]$ irred.\ with same degree $\Rightarrow$ $f$ irred.\ over $\Q$.

\subsection*{11. Euler \& Number Theory}

\textbf{Euler totient.} $\varphi(n)$ = count of $1 \le a \le n$ with $\gcd(a,n) = 1$.

\textbf{Formulas:}
\begin{itemize}[leftmargin=*, itemsep=0pt, topsep=2pt]
\item $\varphi(p^k) = p^{k-1}(p-1)$
\item $\varphi(mn) = \varphi(m)\varphi(n)$ when $\gcd(m,n) = 1$
\item $\varphi(n) = n \prod_{p \mid n}(1 - 1/p)$
\end{itemize}

\textbf{Product of primes:}
\begin{itemize}[leftmargin=*, itemsep=0pt, topsep=2pt]
\item Distinct: $\varphi(pqr) = (p{-}1)(q{-}1)(r{-}1)$
\item Two equal: $\varphi(p^2 q) = p(p{-}1)(q{-}1)$
\item All equal: $\varphi(p^3) = p^2(p{-}1)$
\end{itemize}

\textbf{Units.} $(\Z/n\Z)^\times$ has order $\varphi(n)$.

\textbf{Euler's Thm.} $\gcd(a,n) = 1 \Rightarrow a^{\varphi(n)} \equiv 1 \pmod{n}$.

\textbf{Fermat.} $p$ prime $\Rightarrow$ $a^{p-1} \equiv 1 \pmod{p}$ (if $p \nmid a$).

\subsection*{12. Chinese Remainder Theorem}

\textbf{CRT isomorphism.} $\gcd(m,n) = 1 \Rightarrow \Z/mn\Z \cong \Z/m\Z \times \Z/n\Z$.

\textbf{Converse.} $\gcd(m,n) > 1 \Rightarrow \Z/mn\Z \not\cong \Z/m\Z \times \Z/n\Z$.

\textit{Why?} $\chr(\Z/m\Z \times \Z/n\Z) = \lcm(m,n) \neq mn$.

\textbf{Unit count.} $|(\Z/m\Z \times \Z/n\Z)^\times| = \varphi(m)\varphi(n)$.

\textbf{Solving systems.} $x \equiv a_i \pmod{n_i}$ with coprime $n_i$: unique solution mod $N = \prod n_i$.

\textbf{Non-coprime.} Solvable iff $a_i \equiv a_j \pmod{\gcd(n_i, n_j)}$ for all $i,j$.

\subsection*{13. Ring Isomorphism Invariants}

To show $A \not\cong B$, find an invariant that differs:
\begin{itemize}[leftmargin=*, itemsep=0pt, topsep=2pt]
\item Characteristic: $\chr(A) \neq \chr(B)$
\item Size: $|A| \neq |B|$
\item Unit count: $|A^\times| \neq |B^\times|$
\item Zero divisor count
\end{itemize}

\textit{Ex: $\chr(\Z/9\Z \times \Z/12\Z) = 36 \neq 108 = \chr(\Z/108\Z)$.}

\end{multicols}

%=== PAGE 2: PROBLEM-SOLVING ===%
\newpage
\begin{multicols}{2}
\raggedcolumns

\section*{Problem-Solving Strategies}

\subsection*{Groups \& Permutations}

\vspace{3pt}
\fbox{\parbox{\dimexpr\linewidth-2\fboxsep-2\fboxrule}{%
\textbf{Prove $G$ is abelian}}}
\vspace{3pt}

\textbf{If $|G| = p^2$:} Always abelian (class equation argument).

\textbf{If $|G| = 4$:} Either cyclic, or all non-$e$ have order 2.

In latter case: $(xy)^2 = e \Rightarrow xy = y^{-1}x^{-1} = yx$.

\textbf{General:} Cyclic $\Rightarrow$ abelian. Or check $xy = yx$ for generators.

\vspace{6pt}
\fbox{\parbox{\dimexpr\linewidth-2\fboxsep-2\fboxrule}{%
\textbf{Classify abelian groups of order $n$}}}
\vspace{3pt}

\begin{enumerate}[leftmargin=*, itemsep=1pt, topsep=2pt]
\item Factor $n = p_1^{a_1} \cdots p_r^{a_r}$.
\item For each $p_i^{a_i}$, list partitions of $a_i$.
\item Total count = $\prod_i p(a_i)$.
\item Write groups using elementary or invariant factors.
\end{enumerate}

\textit{For smallest max order: use partition with smallest largest part.}

\vspace{6pt}
\fbox{\parbox{\dimexpr\linewidth-2\fboxsep-2\fboxrule}{%
\textbf{Show no subgroup of order $k$ exists}}}
\vspace{3pt}

For normal subgroups: must be union of conjugacy classes.

\begin{enumerate}[leftmargin=*, itemsep=1pt, topsep=2pt]
\item List conjugacy class sizes (include 1 for identity).
\item Check if any subset sums to $k$.
\item No valid subset $\Rightarrow$ no such subgroup.
\end{enumerate}

\textit{Ex: $A_4$ classes: 1, 3, 4, 4. No subset sums to 6.}

\vspace{6pt}
\fbox{\parbox{\dimexpr\linewidth-2\fboxsep-2\fboxrule}{%
\textbf{Simplify non-disjoint cycle products}}}
\vspace{3pt}

$(1\,2\cdots k) \cdot (i\,j)$: Apply right-to-left. Trace each element.

\textbf{Case 1:} $i \le k < j$ — absorbs $j$ into cycle, gives $(k{+}1)$-cycle.

\textit{Ex: $(1234567)(3\,8) = (1\,2\,3\,8\,4\,5\,6\,7)$.}

\textbf{Case 2:} $i < j \le k$ — splits cycle into two.

\textbf{Case 3:} $k < i < j$ — disjoint, order = $\lcm(k, 2)$.

\vspace{6pt}
\fbox{\parbox{\dimexpr\linewidth-2\fboxsep-2\fboxrule}{%
\textbf{Show $Z(G)$ and $C_G(x)$ are subgroups}}}
\vspace{3pt}

Check subgroup criteria: (1) contains $e$, (2) closed under product, (3) closed under inverse.

For $Z(G)$: if $z_1g = gz_1$ and $z_2g = gz_2$ for all $g$, then $(z_1z_2)g = g(z_1z_2)$.

Note: $C_G(x)$ = stabilizer of $x$ under conjugation action.

\subsection*{Rings \& Polynomials}

\vspace{3pt}
\fbox{\parbox{\dimexpr\linewidth-2\fboxsep-2\fboxrule}{%
\textbf{Analyze quotient ring $F[X]/(p)$}}}
\vspace{3pt}

\begin{itemize}[leftmargin=*, itemsep=1pt, topsep=2pt]
\item \textbf{Size:} $|F|^{\deg p}$
\item \textbf{Elements:} Polynomials of degree $< \deg p$
\item \textbf{Field?} Yes iff $p$ irreducible
\item \textbf{Units:} $\bar{f}$ unit iff $\gcd(f,p) = 1$
\item \textbf{Zero divisors:} $\bar{f}$ ZD iff $\gcd(f,p) \neq 1$
\end{itemize}

\textbf{If field:} Unit group is cyclic of order $|F|^{\deg p} - 1$.

\textit{Ex: $\F_3[X]/(X^2{-}X{-}1)$ is field with 9 elements. Units $\cong C_8$.}

\vspace{6pt}
\fbox{\parbox{\dimexpr\linewidth-2\fboxsep-2\fboxrule}{%
\textbf{Find $\gcd$ in $\F_p[X]$ + Extended Euclidean}}}
\vspace{3pt}

\begin{enumerate}[leftmargin=*, itemsep=1pt, topsep=2pt]
\item Divide: $f = qg + r$, $\deg r < \deg g$. Reduce coeffs mod $p$.
\item Replace: $f \leftarrow g$, $g \leftarrow r$. Repeat until $r = 0$.
\item Last nonzero remainder = $\gcd$. Make monic.
\end{enumerate}

\textbf{Extended:} Track $r, s$ with $\gcd = fr + gs$ (Bézout).

\textit{In $\F_5$: $-1 \equiv 4$, $-2 \equiv 3$, $2^{-1} \equiv 3$, $3^{-1} \equiv 2$.}

\vspace{6pt}
\fbox{\parbox{\dimexpr\linewidth-2\fboxsep-2\fboxrule}{%
\textbf{Test irreducibility}}}
\vspace{3pt}

\textbf{Degree $\le 3$:} No roots $\Rightarrow$ irreducible.

\textbf{Over $\Q$:}
\begin{itemize}[leftmargin=*, itemsep=1pt, topsep=2pt]
\item Rational root test: $r/s$ root $\Rightarrow r \mid a_0$, $s \mid a_n$.
\item Eisenstein with prime $p$.
\item Reduction mod $p$ (same degree, irreducible mod $p$).
\end{itemize}

\textbf{Over $\F_p$:} Test all roots $0, 1, \ldots, p-1$.

\subsection*{Number Theory}

\vspace{3pt}
\fbox{\parbox{\dimexpr\linewidth-2\fboxsep-2\fboxrule}{%
\textbf{Prove $n \mid a$ using modular arithmetic}}}
\vspace{3pt}

\begin{enumerate}[leftmargin=*, itemsep=1pt, topsep=2pt]
\item Factor $n = \prod p_i^{k_i}$.
\item Show $a \equiv 0 \pmod{p_i^{k_i}}$ for each prime power.
\item By CRT, this implies $n \mid a$.
\end{enumerate}

Use Fermat/Euler to reduce exponents.

\textit{Ex: $30 \mid (p^{88} - p^{16})$ for $p \ge 7$? Check mod 2, 3, 5.}

\vspace{6pt}
\fbox{\parbox{\dimexpr\linewidth-2\fboxsep-2\fboxrule}{%
\textbf{Solve CRT system}}}
\vspace{3pt}

Given $x \equiv a_i \pmod{n_i}$ with pairwise coprime $n_i$:

\begin{enumerate}[leftmargin=*, itemsep=1pt, topsep=2pt]
\item Write $x = n_1 t + a_1$.
\item Substitute into next congruence, solve for $t$.
\item Repeat. Solution unique mod $N = \prod n_i$.
\end{enumerate}

Infinitely many solutions: $x = x_0 + Nk$ for $k \in \Z$.

\vspace{6pt}
\fbox{\parbox{\dimexpr\linewidth-2\fboxsep-2\fboxrule}{%
\textbf{Show rings not isomorphic}}}
\vspace{3pt}

Find invariant that differs:
\begin{itemize}[leftmargin=*, itemsep=1pt, topsep=2pt]
\item Characteristic
\item Size
\item Number of units
\item Number of zero divisors
\end{itemize}

\textit{Ex: $\Z/9\Z \times \Z/12\Z$ vs $\Z/108\Z$:}

$\chr = \lcm(9,12) = 36 \neq 108$. Not isomorphic.

$|$units$| = \varphi(9)\varphi(12) = 24 \neq 36 = \varphi(108)$.

\vspace{6pt}
\fbox{\parbox{\dimexpr\linewidth-2\fboxsep-2\fboxrule}{%
\textbf{When does hom exist?}}}
\vspace{3pt}

\textbf{Group hom $C_m \to C_n$:}

Nontrivial exists iff $\gcd(m,n) > 1$.

\textit{Ex: $C_{15} \to C_{30}$: True since $\gcd(15,30)=15 > 1$.}

\textbf{Ring hom $\Z/m\Z \to \Z/n\Z$ (unital):}

Exists iff $n \mid m$ (target divides source).

\textit{Ex: $\Z/15\Z \to \Z/30\Z$: False since $30 \nmid 15$.}

\textit{Why?} $\phi(1) = 1 \Rightarrow \phi(m) = m \cdot 1 = 0 \text{ in target} \Rightarrow n \mid m$.

\subsection*{Quick Reference}

\textbf{Partitions:} $p(1)=1$, $p(2)=2$, $p(3)=3$, $p(4)=5$, $p(5)=7$, $p(6)=11$.

\textbf{Totients:} $\varphi(1)=1$, $\varphi(2)=1$, $\varphi(3)=2$, $\varphi(4)=2$, $\varphi(5)=4$, $\varphi(6)=2$, $\varphi(8)=4$, $\varphi(9)=6$, $\varphi(10)=4$, $\varphi(12)=4$.

\textbf{Common mistakes:}
\begin{itemize}[leftmargin=*, itemsep=0pt, topsep=2pt]
\item $\Z[X]$ is NOT a PID: $(2, X)$ not principal.
\item $\Z/n\Z$ field iff $n$ prime.
\item Ring hom must send $1 \mapsto 1$.
\end{itemize}

\end{multicols}
\end{document}
